% =====================================================================
% §1 INTRODUCTION — reframed draft for real-estate/finance audiences
% Follows the 14-paragraph funnel in
%   skills/realestate-finance-writing/references/introduction.md
% Companion change log: 01-introduction-changes.md
% =====================================================================

\section{Introduction}
\label{sec:intro}

% ---------- ¶1  puzzle, then the theoretical benchmark ----------
Mandatory disclosure is the standard regulatory response to asymmetric information about
product quality. Yet the case for mandating disclosure is contested, because theory suggests
markets may not need it. When quality is verifiable, the unraveling result of
\citet{grossman_informational_1981} and \citet{milgrom_good_1981} holds that sellers disclose
voluntarily, since buyers rationally read silence as the worst possible news. All but the very
worst sellers therefore disclose, and disclosure becomes effectively universal without a
mandate. Whether this result describes real markets remains unsettled. It rests on strong
assumptions, and it is difficult to test, because in most settings where quality is verifiable
disclosure has long been compulsory.

% ---------- ¶2  "to learn about X, we examine Y" ----------
To examine whether unraveling describes a market where quality is verifiable but silence is
common, we study the French residential housing market around the July~2021 enforcement of
Energy Performance Certificate (EPC) disclosure. Under a 2018 reform that took effect on
July~1, 2021, a property's \emph{diagnostic de performance \'energ\'etique} (DPE) became
legally opposable.\footnote{Article 179 III of \emph{loi n\textdegree{}2018-1021 du 23 novembre
2018} (loi \'ELAN) amended Articles L271-4 and L134-3-1 of the \emph{Code de la construction et
de l'habitation}, with opposability effective July~1, 2021; the accompanying methodology
overhaul (replacing the earlier invoice-based calculation with the current
physical-characteristics-based method) was set by the \emph{Arr\^et\'e du 31 mars 2021}. The
\emph{loi Climat et R\'esilience} (\emph{loi n\textdegree{}2021-1104 du 22 ao\^ut 2021}) is a
separate, later reform; we invoke it elsewhere in the paper for the F/G rental-ban calendar it
introduced, not for DPE opposability.} A seller who advertises an energy score is now bound by
it, and buyers acquired legal recourse against misstatement. The reform changed what a disclosed
score means without changing the underlying housing stock or the legal consequences of remaining
silent. It therefore isolates the demand-side value of the signal, and it does so in a market
where the certificate is issued and archived by a national agency rather than by the seller or
the listing platform.

% ---------- ¶3  what theory predicts here, and why this market adjudicates ----------
Applied to this setting, unraveling makes a sharp prediction. Once the DPE is credible, sellers
of energy-efficient properties should publish their score to separate themselves; sellers of
mediocre properties should then publish to avoid being pooled with worse ones; and the process
should continue until listings without a score are confined to the very bottom of the quality
distribution, priced at or below the worst disclosed grade. French housing is a suitable
testbed for this prediction because the disclosed attribute is a single certified index with a
clear mapping into a running cost buyers face directly, because the certificate is verifiable
against an independent public registry, and because a large stock of properties entered the
post-reform market without one.

% ---------- ¶4  the assumption audit ----------
The unraveling prediction rests on four assumptions: (a)~buyers know sellers possess private
information about quality, (b)~disclosure is costless, (c)~disclosure is credible, and
(d)~buyers are homogeneous in how much they value the disclosed attribute. Three of the four are
either directly satisfied in this setting or can be dispensed with. Assumption~(c) holds by
construction after July~2021, since the DPE is issued by a certified third-party diagnostician,
archived in the registry of the \emph{Agence de la Transition \'Ecologique} (ADEME), and legally
binding on the seller. Assumption~(b) is close to satisfied, because a DPE is a one-off outlay of
roughly 100--250~\euro{} according to ADEME's own estimates, negligible relative to
residential transaction values, and it is the seller's obligation rather than the buyer's.
Assumption~(a), which the theoretical literature treats as a maintained hypothesis, we do not
need to assume at all. By matching listings to the ADEME registry we observe directly which
sellers hold a certificate at the moment of listing, and among the apartment listings we can
match at the address level, 48.8\% hold none.
Silence in this market is therefore not always a choice, a possibility first identified by
\citet{dye_disclosure_1985} and formalised by \citet{jung_kwon_1988}, and here it is measured
rather than inferred.

% ---------- ¶5  isolate the assumption at stake ----------
The descriptiveness of unraveling in this market accordingly hinges on assumption~(d).
Unlike disclosure costs and disclosure credibility, which a large literature has examined,
the homogeneity of buyer valuations has received little attention, even though it sets the size of
the disclosure incentive. A seller's penalty for withholding is proportional to the share of
buyers who price the withheld attribute. Housing makes the
assumption particularly implausible. A buyer of a rarely-occupied second home in a mild climate
may rationally ignore an energy score that a buyer of a primary residence in a cold region
prices directly, as may an investor facing the impending rental prohibition on the worst-rated
properties. Evidence across many regulated markets shows that disclosure interventions fail
precisely when recipients are inattentive or heterogeneous in their response
\citep{loewenstein_disclosure_2014, benshahar_failure_2011}, and
\citet{fishman_mandatory_2003} show formally that mandated disclosure can be ineffective, or
welfare-reducing, when buyers differ in their capacity to use the disclosed information. We
therefore let the share of buyers who price the EPC vary across local markets and estimate it.

% ---------- ¶6  data and vantage point  [PRESERVED from current ¶6] ----------
Testing a theory of strategic concealment requires observing both sides of the information
asymmetry at once, the seller's true information, the seller's disclosure choice and the
market's response. Most empirical work observes only one or two. Studies of information
asymmetry in housing infer agent advantages from equilibrium price deviations
\citep{levitt_syverson_2008}; analyses of certification intermediaries characterise incentives
without observing individual information endowments \citep{lizzeri_1999}; and work on the
energy-efficiency gap proxies for seller awareness through indirect variation in energy prices
\citep{myers_asymmetric_2020}. We observe all three. Listing records from SeLoger, a major
French residential platform, give the listing price, property characteristics, the seller's
disclosure decision, and, when a score is displayed, the letter itself. The registry of the
\emph{Agence de la Transition \'Ecologique} (ADEME) is the official national database of
certified DPE audits, reported by accredited third-party assessors and compiled independently
of the listing market. We match listings to the registry on address, property type, area, and
floor, which identifies whether a certified audit existed at the time of listing, the one
thing buyers can never observe. Because the match requires a geocoded address, the matched
sample is the subset of listings that can be located at the address level rather than the full
listing universe. Within that matched sample, every listing falls into one of four groups, namely
truthful disclosers, strategic hiders who hold a certificate and omit it, genuinely uninformed
sellers for whom no certificate exists, and a small residual displaying a score without a
verified audit. The information structure of the disclosure game is therefore observed rather
than inferred, a vantage point that to our knowledge is unique in this literature. Transaction
records enter separately and are not matched to individual listings. We use them to estimate
the mapping from listing to realised sale price, which supports treating the listing price as
the relevant outcome, since the negotiation margin in French residential sales is small and
most properties transact at or near list price
\citep{merlo_bargaining_2004, genesove_loss_2001, friggit_marche_2013}.

% ---------- ¶7  descriptive findings, enumerated ----------
The data yield four descriptive findings, drawn from two distinct samples that we keep separate
throughout. First, silence commands a premium rather than a penalty. Estimated on the panel of
market-level hedonic price indices, and relative to the pre-reform baseline, listings without a
score trade approximately 6.9\% above disclosed energy-intensive properties for apartments and
3.4\% for houses, the opposite of the sign unraveling predicts. Second, the premium widened
abruptly at the reform rather than emerging from trend. The composition-adjusted
difference-in-differences estimates evolve in parallel before July~2021, averaging $+0.002$ for
apartments and $+0.008$ for houses with no discernible drift, and then shift sharply upward,
stabilising around $+0.044$ and $+0.037$ by mid-2024. For apartments the level ordering already
favoured non-disclosers beforehand, so what the reform changes is the magnitude of the premium
rather than its sign. Third, turning to the address-matched listing sample, non-disclosure is
overwhelmingly involuntary. Among apartment listings 79.8\% disclose truthfully, 16.7\% are
genuinely uninformed, and only 3.5\% hold a certificate and withhold it, with the corresponding
house shares 90.4\%, 8.5\% and 1.1\%.
Fourth, the non-disclosure rate varies widely across local markets, ranging among
apartments from 6.5\% to 40.9\% across d\'epartements, around a median of 16.5\%. Together
these facts are inconsistent with unraveling as a description of this market, while pointing to
the mechanism that defeats it. The pool of silent sellers is dominated by the uninformed, which
blunts the inference buyers can draw from silence and shelters the few who conceal
deliberately.

% ---------- ¶8  model motivation ----------
To recover the forces behind these patterns, we build and estimate a model of the disclosure
game that allows for, but does not impose, strategic concealment. Sellers differ in whether they
hold a certified audit and in the quality it would reveal; buyers differ in whether they price
energy performance. Informed sellers of poor properties choose between publishing a bad score
and accepting the price buyers assign to silence, and that choice depends on how contaminated
the silent pool is by sellers who simply have no certificate. Because the composition of the
silent pool is itself the outcome of those choices, the model is closed by a fixed point in
buyer beliefs, and it nests the classical benchmark. When audit coverage and buyer sensitivity
are both high, the silent pool empties and the Grossman--Milgrom outcome is recovered. For
intermediate values, which is where the French data sit, the model predicts partial pooling in
which some informed sellers conceal and silence carries a premium.

% ---------- ¶9  what the model finds ----------
The estimates deliver two structural magnitudes and one specification test. Identifying buyer
energy-sensitivity separately from the concealment cost requires a source of variation external
to prices, since the two are shown to be identified only as their product by price moments
alone (Appendix~\ref{app:rho_identification}); we construct one from the capitalised energy-cost
saving a sensitive buyer attaches to verified high- versus low-quality certification, built from
ADEME-certified consumption bands rather than from the price data itself. The resulting share of
buyers who price energy performance averages 49.5\% across apartment markets and 45.8\% across
house markets, and the cost of withholding a certificate settles at 2.4\% of property value for
apartments and 3.5\% for houses, on the order of a third of the notary bill a French buyer
separately pays on the same transaction. Both are economically substantial magnitudes, and of the
same order in two segments that trade very differently. The model predicts, in the partial-pooling
equilibrium that describes the bulk of French housing markets, that the non-disclosure premium
should collapse to this concealment cost exactly, invariant to market composition; we test this
prediction directly rather than impose it. It holds for houses: the gap between the observed
premium and the identified concealment cost is small and statistically indistinguishable from
zero. It does not hold for apartments: the premium exceeds the concealment cost by a wide and
precisely estimated margin, a specification tension the baseline cannot absorb, because no belief
about silence can price a silent listing above a verified efficient one.
% [AUTHOR REWRITE]
Section~\ref{ss:attention} shows that one change on the buyer side can account for it: if buyers attend
to energy when a rating is displayed and much less when it is not, the model fits houses as before,
rejects full attention in apartments, and leaves the estimated buyer sensitivity unchanged.
We maintain, rather than test, the model's prediction that informed sellers of high-quality
properties never conceal; Appendix~\ref{app:robustness} relaxes it by reading the withheld grade
from the registry and finds the structural estimates essentially unchanged. Under the limited-attention
extension the prediction survives whenever not every informed low-quality seller hides (Corollary~\ref{cor:att_hq}).

% ---------- ¶11  contribution 1 ----------
Our paper contributes to the literature on disclosure by testing the unraveling result in a
market where the informed and uninformed can be told apart. Prior work has examined the
consequences of uncertainty about information endowment \citep{dye_disclosure_1985,
jung_kwon_1988}, of disclosure costs, and of disclosure credibility, but has generally had to
treat sellers' information as latent. By observing it directly we can decompose silence into
inability and choice, and show that the first dominates the second by an order of magnitude
while the second is what sustains the price premium. The finding that silence is rewarded rather
than punished is not a failure of the theory so much as a demonstration of how sensitive its
central prediction is to an assumption that is rarely examined.

% ---------- ¶12  contribution 2 + scope disclaimer ----------
Our paper also provides an information-based account of the residential energy-efficiency gap
\citep{allcott_is_2012, gerarden_assessing_2017}. The standard diagnosis appeals to split
incentives, credit constraints, and inattention to future energy costs \citep{allcott_consumers_2011, sallee_rational_2014, allcott_evaluating_2015}. We identify a
complementary channel internal to the disclosure game. A seller holding a poor rating can
renovate and publish a good one, earning the quality premium, or stay silent and earn the
non-disclosure premium instead, so the premium to opacity looks like an opportunity cost of
renovation that raises the private hurdle above engineering cost. In equilibrium it does not, except in the corner where every informed low-quality seller hides. This has an unexpected policy
implication. Because the equilibrium absorbs enforcement through the fraction of sellers who
conceal, the renovation threshold depends only on the price gap between disclosed quality
levels and is independent of enforcement intensity. Renovation subsidies, enforcement against
concealment, and investment in audit coverage act on three distinct margins and do not substitute for one another,
contrary to the conventional treatment in
\citet{gillingham_palmer_2014}. The result holds outside the corner in which every informed low-quality seller hides, and with one concealment cost per market. % [AUTHOR REWRITE] The extension of Section~\ref{ss:attention} brings the inattention
diagnosis inside the disclosure game: buyers who attend less to energy when no rating is shown let
silence earn a premium that no belief about quality can produce. We do not claim to have estimated the one true model of housing
disclosure; we propose one way of reading these data through a workhorse framework, and we are
explicit throughout about which parameters are point-identified and which are bounded.

% ---------- ¶12.5  policy extension results ----------
The model also speaks to renovation policy. Extending the framework to an endogenous upgrade
decision, we show the minimum subsidy needed to make renovation privately worthwhile reduces, in
the empirically dominant partial-pooling regime, to the price gap between verified high and low
quality alone, invariant to enforcement against hiding and to how many sellers are informed,
which act on separate margins instead. Forcing silent listings to be read as carefully as rated ones would change who hides but not who renovates, and only a higher weight on displayed ratings could raise the renovation threshold, within a bound that the data do not pin down. % [AUTHOR REWRITE]
A uniform national subsidy is accordingly poorly targeted. The
required cost-share the state must cover varies widely across
markets, averaging \attBaseCostShareTwoFiftyApt{} percent in apartment markets that need support at all at the light benchmark, and aggregate programme costs among the sellers who can renovate run to
\euro{}\attScCostBaseTotTwoFifty{}--\attScCostBaseTotFiveHundred{} billion a year in our listing sample at plausible renovation-cost benchmarks, for scale between \attScCostBaseTotTwoFiftyLo{} and
\attScCostBaseTotFiveHundredHi{} percent of MaPrimeR\'enov's current annual budget. Efficient targeting requires conditioning subsidy design
on local market composition rather than applying one national rate. Enforcement against strategic hiding and investment in audit
infrastructure act on separate margins: the first changes who hides, the second who is eligible to renovate, and neither lowers the subsidy a given seller needs.

% ---------- ¶13  related literature, three strands ----------
Our paper draws on three strands of work. It belongs to the literature that estimates
disclosure models structurally in order to assess the descriptiveness of classical disclosure
theory and to recover primitives such as information endowment, disclosure costs, and beliefs;
closest to us is \citet{bourveau_learning_2025}, who structurally examine whether the unraveling
result describes an unregulated nineteenth-century market and find that full disclosure is
reached only gradually. Where that literature and the experimental work on disclosure ask how far
receivers infer the worst from silence \citep{jin_is_2021, hagenbach_motivated_2025}, our extension
asks whether they attend to silence at all, drawing on evidence that limited attention changes how
disclosed and undisclosed information is priced \citep{hirshleifer_limited_2003, dellavigna_investor_2009}.
It is closely related to the extensive evidence on the capitalisation of energy labels in
housing, which documents robust premia for efficient properties wherever labels have been
studied \citep{brounen_kok_2011, hyland_lyons_2013, fuerst_energy_2016, frondel_power_2018,
myers_mandatory_2022}; that literature establishes that labels are priced, whereas our question
is why their \emph{absence} is priced favourably. And it builds on work on information
asymmetry in real-estate markets, where the informational advantage of one side of the
transaction has been inferred from prices and time-on-market \citep{levitt_syverson_2008,
lizzeri_1999}.

% ---------- ¶14  generalizability, conceded then defended ----------
The features that make this setting informative also bound it. French EPC disclosure is a
single mandate, in one country, covering one verifiable attribute, and the low audit coverage
that lets us separate inability from choice is itself a feature of a young disclosure regime.
The mechanism, however, is not specific to energy labels. Wherever a disclosure mandate is
introduced before the infrastructure that makes disclosure feasible is universal, a pool of
genuinely uninformed non-disclosers will exist, and that pool will shelter strategic
concealment and dull the incentive to disclose. That configuration characterises a broad class
of emerging disclosure regimes, from climate-risk and emissions reporting to supply-chain and
sustainability disclosure, in which mandates routinely precede measurement capacity, including
the French carbon-disclosure mandate studied by \citet{bourveau_carbon_2024}, where the same
gap between obligation and measurement infrastructure arises.
